\Needspace{8\baselineskip}
\Article{191}

\textsc{Satz.} \textit{Geht eine reduzierte Form $(a,b,-a')$ mit Determinante $D$ durch die Substitution $\alpha,\beta,\gamma,\delta$ in die reduzierte Form $(A,B,-A')$ mit derselben Determinante über, so liegt} erstens \textit{$\frac{\pm\surd D-b}{a}$ zwischen $\frac{\alpha}{\gamma}$ und $\frac{\beta}{\delta}$ (sofern weder $\gamma$ noch $\delta=0$ ist, also beide Grenzen endlich sind), wobei das obere Vorzeichen zu nehmen ist, wenn keine der beiden Grenzen das zu $a$ entgegengesetzte Vorzeichen hat (oder deutlicher: wenn entweder beide das Vorzeichen dieses Koeffizienten haben oder eine dieses Vorzeichen hat und die andere $=0$ ist), das untere, wenn keine dasselbe Vorzeichen hat wie $a$;} zweitens \textit{$\frac{\pm\surd D+b}{a'}$ zwischen $\frac{\gamma}{\alpha}$ und $\frac{\delta}{\beta}$ (sofern weder $\alpha$ noch $\beta=0$ ist), wobei das obere Vorzeichen zu nehmen ist, wenn keine Grenze das zu $a'$ (oder $a$) entgegengesetzte Vorzeichen hat, das untere, wenn keine dasselbe Vorzeichen hat wie $a'$}\HistoricalNote{*)}{Andere Fälle können offenbar nicht auftreten, da nach dem vorigen Artikel wegen $\alpha\delta-\beta\gamma=\pm1$ die beiden Grenzen weder entgegengesetzte Vorzeichen haben noch beide zugleich $=0$ sein können.}.

\textit{Beweis.} Wir haben die Gleichungen
\[
\begin{array}{rcll}
a\alpha\alpha+2b\alpha\gamma-a'\gamma\gamma&=&A&\qquad\mathinner{\ldotp\ldotp\ldotp\ldotp\ldotp\ldotp\ldotp\ldotp}\qquad[1]\\[4pt]
a\beta\beta+2b\beta\delta-a'\delta\delta&=&-A'&\qquad\mathinner{\ldotp\ldotp\ldotp\ldotp\ldotp\ldotp\ldotp}\qquad[2]
\end{array}
\]
Daraus folgt
\[
\begin{array}{rcll}
\dfrac{\alpha}{\gamma}&=&\dfrac{\pm\surd\bigl(D+\frac{aA}{\gamma\gamma}\bigr)-b}{{a}}&\qquad\mathinner{\ldotp\ldotp\ldotp\ldotp\ldotp\ldotp\ldotp\ldotp}\qquad[3]\\[12pt]
\dfrac{\beta}{\delta}&=&\dfrac{\pm\surd\bigl(D-\frac{aA'}{\delta\delta}\bigr)-b}{a}&\qquad\mathinner{\ldotp\ldotp\ldotp\ldotp\ldotp\ldotp\ldotp\ldotp}\qquad[4]
\end{array}
\]
\SourceLeaf{182}
\begin{equation*}
\frac{\gamma}{\alpha}=\frac{\pm\surd\bigl(D-\frac{a'A}{\alpha\alpha}\bigr)+b}{a'}
\quad \mathinner{\ldotp\ldotp\ldotp\ldotp\ldotp\ldotp\ldotp\ldotp} \tag*{[5]}
\end{equation*}
\begin{equation*}
\frac{\delta}{\beta}=\frac{\pm\surd\bigl(D+\frac{a'A'}{\beta\beta}\bigr)+b}{a'}
\quad \mathinner{\ldotp\ldotp\ldotp\ldotp\ldotp\ldotp\ldotp\ldotp} \tag*{[6]}
\end{equation*}
Die Gleichung 3, 4. 5, 6 ist jeweils wegzulassen, wenn $\gamma,\delta,\alpha,\beta$ entsprechend $=0$ ist. --- Es bleibt aber offen, welche \textit{Vorzeichen} den Wurzelgrößen zu geben sind. Dies entscheiden wir so.

In [3] und [4] müssen offenbar die oberen Vorzeichen genommen werden, wenn weder $\frac{\alpha}{\gamma}$ noch $\frac{\beta}{\delta}$ das zu $a$ entgegengesetzte Vorzeichen hat; denn beim unteren Vorzeichen würden $\frac{a\alpha}{\gamma}$ und $\frac{a\beta}{\delta}$ negative Größen. Da aber $A$ und $A'$ dasselbe Vorzeichen haben, liegt $\surd D$ zwischen $\surd\bigl(D+\frac{aA}{\gamma\gamma}\bigr)$ und $\surd\bigl(D-\frac{aA'}{\delta\delta}\bigr)$ und daher in diesem Fall $\frac{\surd D-b}{a}$ zwischen $\frac{\alpha}{\gamma}$ und $\frac{\beta}{\delta}$. Damit ist der erste Teil des Satzes für den ersten Fall bewiesen.

Ebenso sieht man, dass in [5] und [6] die unteren Vorzeichen genommen werden müssen, wenn weder $\frac{\gamma}{\alpha}$ noch $\frac{\delta}{\beta}$ dasselbe Vorzeichen hat wie $a'$ oder $a$, denn beim oberen Vorzeichen würden $\frac{a'\gamma}{\alpha}$, $\frac{a'\delta}{\beta}$ notwendig positive Größen. Daraus folgt sofort, dass $\frac{-\surd D+b}{a'}$ in diesem Fall zwischen $\frac{\gamma}{\alpha}$ und $\frac{\delta}{\beta}$ liegt. Damit ist auch der zweite Teil des Satzes für den zweiten Fall bewiesen. Ließe sich ebenso leicht zeigen, dass in [3] und [4] die unteren Vorzeichen zu nehmen sind, wenn keine der Größen $\frac{\alpha}{\gamma}$, $\frac{\beta}{\delta}$ dasselbe Vorzeichen hat wie $a$, und in [5] und [6] die oberen, wenn weder $\frac{\gamma}{\alpha}$ noch $\frac{\delta}{\beta}$ das entgegengesetzte Vorzeichen hat, so würde ebenso folgen, dass im ersteren Fall $\frac{-\surd D-b}{a}$ zwischen $\frac{\alpha}{\gamma}$ und $\frac{\beta}{\delta}$ liegt, im letzteren $\frac{\surd D+b}{a'}$ zwischen $\frac{\gamma}{\alpha}$ und $\frac{\delta}{\beta}$. Dann wären der erste Teil auch für den zweiten Fall und der zweite Teil für den ersten Fall bewiesen. Obwohl das nicht schwierig ist, geht es aber nicht ohne einige Umwege. Wir ziehen die folgende Methode vor.

Ist keine der Zahlen $\alpha,\beta.\gamma,\delta=0$, so haben $\frac{\alpha}{\gamma}$ und $\frac{\beta}{\delta}$ dieselben Vorzeichen wie $\frac{\gamma}{\alpha}$, $\frac{\delta}{\beta}$. Hat also keine dieser Größen dasselbe Vorzeichen wie $a'$ oder $a$, sodass $\frac{-\surd D+b}{a'}$ zwischen $\frac{\gamma}{\alpha}$ und $\frac{\delta}{\beta}$ liegt, so hat keine der Größen $\frac{\alpha}{\gamma}$, $\frac{\beta}{\delta}$ dasselbe Vorzeichen wie $a$, und $\frac{a'}{-\surd D+b}=\frac{-\surd D-b}{a}$ liegt (wegen $aa'=D-bb$) zwischen $\frac{\alpha}{\gamma}$ und $\frac{\beta}{\delta}$. Damit ist, sofern weder $\alpha$ noch $\beta=0$ ist, der erste Teil des Satzes auch für den zweiten Fall bewiesen (denn die Bedingung, dass weder $\gamma$ noch $\delta=0$ ist, steht bereits im Satz selbst). Ebenso hat, wenn keine der Zahlen $\alpha,\beta,\gamma,\delta=0$ ist
\SourceLeaf{183}
und weder $\frac{\alpha}{\gamma}$ noch $\frac{\beta}{\delta}$ das zu $a$ oder $a'$ entgegengesetzte Vorzeichen hat, sodass ${\frac{\surd D-b}{a'}}$ zwischen $\frac{\alpha}{\gamma}$ und $\frac{\beta}{\delta}$ liegt, auch keine der Größen $\frac{\gamma}{\alpha}$ und $\frac{\delta}{\beta}$ das zu $a'$ entgegengesetzte Vorzeichen, und $\frac{a}{\surd D-b}=\frac{\surd D+b}{a'}$ liegt zwischen $\frac{\gamma}{\alpha}$ und $\frac{\delta}{\beta}$. Damit ist für den Fall, dass weder $\gamma$ noch $\delta=0$ ist, der zweite Teil des Satzes auch für den zweiten Fall bewiesen.

Es bliebe also nur zu beweisen, dass der erste Teil des Satzes auch für den zweiten Fall gilt, wenn eine der Zahlen $\alpha,\beta$ $=0$ ist, und der zweite Teil für den ersten Fall, wenn entweder $\gamma$ oder $\delta=0$ ist. Aber \textit{alle diese Fälle sind unmöglich}. Für den \textit{ersten} Teil des Satzes nehmen wir an, weder $\gamma$ noch $\delta=0$ sei, $\frac{\alpha}{\gamma}$, $\frac{\beta}{\delta}$ hätten nicht dasselbe Vorzeichen wie $a$, und es sei 1) $\alpha=0$. Dann folgt aus der Gleichung $\alpha\delta-\beta\gamma=\pm1$, dass $\beta=\pm1$, $\gamma=\pm1$ ist. Aus [1] folgt $A=-a'$, also haben $A$ und $a'$ und damit auch $a$ und $A'$ entgegengesetzte Vorzeichen; daher gilt $\surd\bigl(D-\frac{aA'}{\delta\delta}\bigr)>\surd D>b$. In [4] ist somit das untere Vorzeichen zu nehmen, denn beim oberen hätte $\frac{\beta}{\delta}$ offenbar dasselbe Vorzeichen wie $a$. Also gilt $\frac{\beta}{\delta}>\frac{-\surd D-b}{a}>1$ (weil nach der Definition einer reduzierten Form $a<\surd D+b$ ist), \textit{was widersinnig ist}, da $\beta=\pm1$ und $\delta$ nicht $=0$ ist. --- 2) Sei $\beta=0$. Dann folgt aus der Gleichung $\alpha\delta-\beta\gamma=\pm1$, dass $\alpha=\pm1$, $\delta=\pm1$ ist. Aus [2] folgt $-A'=-a'$, sodass $a'$ und $a$ und $A$ dasselbe Vorzeichen haben; daher gilt ${\surd\bigl(D+\frac{aA}{\alpha\alpha}\bigr)}>\surd D>b$. In [3] ist somit das untere Vorzeichen zu nehmen, denn beim oberen hätte $\frac{\alpha}{\gamma}$ dasselbe Vorzeichen wie $a$. Also gilt $\frac{\alpha}{\gamma}>\frac{-\surd D-b}{a}>1$, \textit{was widersinnig ist}, aus demselben Grund wie zuvor. --- Für den \textit{zweiten} Teil nehmen wir an, weder $\alpha$ noch $\beta=0$ sei, $\frac{\gamma}{\alpha}$, $\frac{\delta}{\beta}$ hätten nicht das zu $a'$ entgegengesetzte Vorzeichen, und es sei 1) $\gamma=0$. Aus der Gleichung $\alpha\delta-\beta\gamma=\pm1$ folgt $\alpha=\pm1$, $\delta=\pm1$. Aus [1] folgt $A=a$, sodass $a'$ und $A'$ dasselbe Vorzeichen haben; daher gilt $\surd\bigl(D+\frac{a'A'}{\beta\beta}\bigr)>\surd D>b$. In [6] ist also das obere Vorzeichen zu nehmen, denn beim unteren hätte $\frac{\delta}{\beta}$ das zu $a'$ entgegengesetzte Vorzeichen. Somit gilt $\frac{\delta}{\beta}>\frac{\surd D+b}{a'}>1$, \textit{was widersinnig ist}, da $\delta=\pm1$ und $\beta$ nicht $=0$ ist. Ist schließlich 2) $\delta=0$, so folgt aus $\alpha\delta-\beta\gamma=\pm1$, dass $\beta=\pm1$, $\gamma=\pm1$ ist, und damit aus [2], dass $-A'=a$. Daraus folgt $\surd\bigl(D-\frac{a'A}{\alpha\alpha}\bigr)>\surd D>b$, weshalb in [5] das obere Vorzeichen zu nehmen ist. Also gilt $\frac{\gamma}{\alpha}>\frac{\surd D+b}{a'}>1$, \textit{was widersinnig ist}. --- \mbox{Damit ist der Satz in seinem ganzen Umfang bewiesen.}

\Needspace{9\baselineskip}
Da die Differenz zwischen $\frac{\alpha}{\gamma}$ und $\frac{\beta}{\delta}$ $=\frac{1}{\gamma\delta}$ ist, ist die Differenz zwischen $\frac{\pm\surd D-b}{a}$ und $\frac{\alpha}{\gamma}$ oder $\frac{\beta}{\delta}$ $<\frac{1}{\gamma\delta}$; und zwischen $\frac{\pm\surd D-b}{a}$ und $\frac{\alpha}{\gamma}$ oder zwischen jener Größe und $\frac{\beta}{\delta}$ kann kein Bruch liegen, dessen Nenner nicht größer als $\gamma$ oder $\delta$ ist (\textit{voriger Hilfssatz}). --- Ebenso ist die Differenz der Größe ${\frac{\pm\surd D+b}{a}}$ zum Bruch $\frac{\gamma}{\alpha}$ oder $\frac{\delta}{\beta}$
\SourceLeaf{184}
kleiner als $\frac{1}{\alpha\beta}$, und zwischen jener Größe und einem dieser Brüche kann kein Bruch liegen, dessen Nenner nicht größer als $\alpha$ und $\beta$ ist.

\Needspace{8\baselineskip}
\Article{192}
Wendet man den vorigen Satz auf den Algorithmus von Artikel~188 an, so folgt, dass die Größe $\frac{\surd D-b}{a}$, die wir mit $L$ bezeichnen, zwischen $\frac{\alpha'}{\gamma'}$ und $\frac{\beta'}{\delta'}$, zwischen $\frac{\alpha''}{\gamma''}$ und $\frac{\beta''}{\delta''}$, zwischen $\frac{\alpha'''}{\gamma'''}$ und $\frac{\beta'''}{\delta'''}$ usw. liegt (denn aus dem Ende von Artikel~189,~3 folgt leicht, dass keine dieser Grenzen das zu $a$ entgegengesetzte Vorzeichen hat, weshalb der Wurzelgröße $\surd D$ das positive Vorzeichen zu geben ist), also zwischen $\frac{\alpha'}{\gamma'}$ und $\frac{\alpha''}{\gamma''}$, zwischen $\frac{\alpha''}{\gamma''}$ und $\frac{\alpha'''}{\gamma'''}$ usw. Daher liegen alle Brüche $\frac{\alpha'}{\gamma'}$, $\frac{\alpha'''}{\gamma'''}$, $\frac{\alpha'''''}{\gamma'''''}$ usw. auf derselben Seite von $L$ und alle $\frac{\alpha''}{\gamma''}$, $\frac{\alpha''''}{\gamma''''}$, $\frac{\alpha''''''}{\gamma''''''}$ usw. auf der anderen. Da aber $\gamma'<\gamma'''$ gilt, liegt $\frac{\alpha'}{\gamma'}$ \textit{außerhalb} des Intervalls zwischen $\frac{\alpha'''}{\gamma'''}$ und $L$, und aus ähnlichem Grund $\frac{\alpha''}{\gamma''}$ außerhalb von $L$ und $\frac{\alpha''''}{\gamma''''}$, $\frac{\alpha'''}{\gamma'''}$ außerhalb von $L$ und $\frac{\alpha'''''}{\gamma'''''}$ usw. Daher liegen diese Größen offenbar in folgender Reihenfolge:
\[
\frac{\alpha'}{\gamma'},\ \frac{\alpha'''}{\gamma'''},\ \frac{\alpha'''''}{\gamma'''''}\mathinner{\ldotp\ldotp\ldotp} L\mathinner{\ldotp\ldotp\ldotp}\frac{\alpha''''''}{\gamma''''''},\ \frac{\alpha''''}{\gamma''''},\ \frac{\alpha''}{\gamma''}
\]
Die Differenz zwischen $\frac{\alpha'}{\gamma'}$ und $L$ ist aber kleiner als die Differenz zwischen $\frac{\alpha'}{\gamma'}$ und $\frac{\alpha''}{\gamma''}$, \textit{das heißt}, $<\frac{1}{\gamma'\gamma''}$, und ebenso ist die Differenz zwischen $\frac{\alpha''}{\gamma''}$ und $L$ $<\frac{1}{\gamma''\gamma'''}$ usw. Daher nähern sich die Brüche $\frac{\alpha'}{\gamma'}$, $\frac{\alpha''}{\gamma''}$, $\frac{\alpha'''}{\gamma'''}$ usw. der Grenze $L$ immer mehr, und da $\gamma',\gamma'',\gamma'''$ fortwährend ins Unendliche wachsen, kann die Differenz der Brüche zur Grenze kleiner als jede gegebene Größe werden.

Nach Artikel~189 hat keine der Größen $\frac{\gamma}{\alpha}$, $\frac{{}'\gamma}{{}'\alpha}$, $\frac{{}''\gamma}{{}''\alpha}$ usw. dasselbe Vorzeichen wie $a$. Durch einen ganz ähnlichen Schluss wie zuvor folgt daraus, dass jene Größen und $\frac{-\surd D+b}{a'}$, die wir mit $L'$ bezeichnen, in folgender Reihenfolge liegen:
\[
\frac{\gamma}{\alpha},\ \frac{{}''\gamma}{{}''\alpha},\ \frac{{}''''\gamma}{{}''''\alpha}\mathinner{\ldotp\ldotp\ldotp\ldotp\ldotp} L'\mathinner{\ldotp\ldotp\ldotp\ldotp\ldotp}\frac{{}'''''\gamma}{{}'''''\alpha},\ \frac{{}'''\gamma}{{}'''\alpha},\ \frac{{}'\gamma}{{}'\alpha}
\]
Die Differenz zwischen $\frac{\gamma}{\alpha}$ und $L'$ ist kleiner als $\frac{1}{{}'\alpha\alpha}$, die Differenz zwischen $\frac{{}'\gamma}{{}'\alpha}$ und $L'$ kleiner als $\frac{1}{{}''\alpha{}'\alpha}$ usw. Daher nähern sich die Brüche $\frac{\gamma}{\alpha}$, $\frac{{}'\gamma}{{}'\alpha}$ usw. immer mehr $L'$, und die Differenz kann kleiner als jede gegebene Größe werden.

Im \textit{Beispiel} von Artikel~188 gilt $L=\frac{\surd79-{8}}{3}=0{,}2960648$, und die Näherungsbrüche sind $\frac{0}{1}$, $\frac{1}{3}$, $\frac{2}{7}$, $\frac{3}{10}$, $\frac{5}{17}$, $\frac{8}{27}$, $\frac{45}{152}$, $\frac{143}{483}$ usw. Dabei ist $\frac{143}{483}=0{,}2960662$. --- Dort gilt auch $L'=\frac{-\surd79+{8}}{5}=-0{,}1776388$, und die Näherungsbrüche sind $\frac{0}{1}$, $-\frac{1}{5}$, $-\frac{1}{6}$, $-\frac{2}{11}$, $-\frac{3}{17}$, $-\frac{8}{45}$, $-\frac{27}{152}$, $-\frac{143}{805}$ usw. Es ist aber $\frac{143}{805}=0{,}1776397$.
